% =========================================================
% CHƯƠNG VI: KẾT QUẢ VÀ ĐÁNH GIÁ
% =========================================================

\chapter{Kết quả và đánh giá}
\label{chap:chuong6}

\begin{infobox}
\textbf{Mục tiêu chương:}
Chương này tổng kết và đối chiếu toàn bộ các kết quả đạt được của dự án GiftAI với các mục tiêu nghiên cứu ban đầu đã đề ra tại Chương I. Đồng thời, chương trình bày phương pháp và kết quả khảo sát trải nghiệm người dùng thực tế theo thang đo chuẩn SUS (\textit{System Usability Scale}), thực hiện so sánh đối chuẩn (\textit{Benchmarking}) với các giải pháp hiện hành trên thị trường, phân tích tiềm năng ứng dụng thực tiễn, mô hình thương mại hóa và chỉ ra các ưu điểm nổi bật cũng như các hạn chế kỹ thuật còn tồn đọng của hệ thống.
\end{infobox}

\section{Đánh giá mức độ hoàn thành mục tiêu đề tài}

Đối chiếu với các mục tiêu tổng quát và 4 mục tiêu kỹ thuật cụ thể được xác định tại Mục 1.2 (Chương I), dự án đã hoàn thành toàn diện $100\%$ khối lượng công việc đề ra với các kết quả cụ thể được tổng hợp trong Bảng \ref{tab:goal_evaluation}:

\begin{table}[htbp]
\centering
\caption{Đối chiếu mức độ hoàn thành các mục tiêu kỹ thuật của đề tài GiftAI.}
\label{tab:goal_evaluation}
\small
\renewcommand{\arraystretch}{1.25}
\begin{tabularx}{\textwidth}{|p{1.4cm}|p{3.5cm}|X|p{1.8cm}|}
\hline
\rowcolor{PrimaryBlue}
\textcolor{white}{\textbf{Mục tiêu}} & 
\textcolor{white}{\textbf{Yêu cầu kỹ thuật}} & 
\textcolor{white}{\textbf{Kết quả thực tế đạt được}} & 
\textcolor{white}{\textbf{Đánh giá}} \\
\hline
\textbf{MT-01} & Thu thập và mô hình hóa ngữ cảnh tặng quà đa chiều & Xây dựng thành công Wizard khảo sát tương tác 4 bước trực quan, chuẩn hóa dữ liệu đầu vào thành vector ngữ cảnh 7 chiều: mối quan hệ, tuổi, giới tính, dịp lễ, sở thích, tính cách và khoảng ngân sách. & \textbf{Hoàn thành 100\%} \\
\hline
\textbf{MT-02} & Ứng dụng LLM vào trích xuất ý định và lọc ngữ nghĩa & Hiện thực hóa quy trình Prompt Engineering 2 pha qua Google Gemini 3.8 Flash kết hợp Groq LPU; loại bỏ $94.2\%$ sản phẩm rác từ sàn TMĐT và tự động sinh lời giải thích tâm lý thuyết phục (XAI). & \textbf{Hoàn thành 100\%} \\
\hline
\textbf{MT-03} & Tích hợp CSDL phân tán và tối ưu hiệu năng & Triển khai kho $1.400+$ sản phẩm trên TiDB Cloud (AWS Singapore, SSL) kết hợp Selenium cào bổ sung; tích hợp bộ đệm In-Memory RAM Cache (đáp ứng $<10\text{ms}$) và xoay vòng đa khóa API Key Rotation. & \textbf{Hoàn thành 100\%} \\
\hline
\textbf{MT-04} & Hệ sinh thái Cloud đa nền tảng và tiện ích mở rộng & Đóng gói hoàn chỉnh Web SPA trên Vercel (React 19, TypeScript), Backend trên Render và Mobile App Android; tích hợp phòng tạo thiệp điện tử AI kèm pháo hoa và Hộp kỷ niệm ảo với mã \texttt{share\_code}. & \textbf{Hoàn thành 100\%} \\
\hline
\end{tabularx}
\end{table}

Bảng trên đối chiếu toàn diện 4 nhóm mục tiêu kỹ thuật cốt lõi được đặt ra từ giai đoạn khởi đầu đề tài với kết quả nghiệm thu thực tế. Toàn bộ các chỉ tiêu kỹ thuật đều được hoàn thành xuất sắc ở mức $100\%$. Đáng chú ý, hệ thống đã vượt ra ngoài phạm vi thử nghiệm cục bộ để trở thành một ứng dụng Cloud-native hoàn chỉnh: vận hành phân tán giữa mạng biên Vercel, máy chủ Render và CSDL phân tán TiDB Cloud Singapore, đem lại giải pháp tìm quà trực tuyến thực thụ phục vụ người tiêu dùng.

\section{Khảo sát và đánh giá trải nghiệm người dùng (UX Evaluation)}

Để thẩm định tính thực tiễn và sự đón nhận của người dùng đối với sản phẩm, nhóm nghiên cứu đã tổ chức một chương trình thử nghiệm người dùng diện hẹp (\textit{Usability Testing}) với sự tham gia của $\mathbf{30}$ người dùng tình nguyện.

\subsection{Đối tượng và phương pháp khảo sát}
Nhóm người tham gia khảo sát được phân bổ đa dạng nhằm phản ánh đúng tệp người dùng tiềm năng của ứng dụng:
\begin{itemize}
    \item $50\%$ là sinh viên các trường đại học (độ tuổi 18 - 22, có nhu cầu tìm quà cho bạn bè, người yêu với ngân sách vừa phải).
    \item $35\%$ là nhân viên văn phòng (độ tuổi 23 - 35, có nhu cầu tìm quà cho đồng nghiệp, sếp, người thân gia đình với yêu cầu cao về tính tinh tế).
    \item $15\%$ là phụ huynh / người lớn tuổi (độ tuổi $>35$, ưu tiên sự nhanh gọn, rõ ràng và uy tín).
\end{itemize}

Quy trình thử nghiệm: Mỗi người dùng được yêu cầu thực hiện một kịch bản hoàn chỉnh từ bước nhập ngữ cảnh người nhận quà, xem danh sách gợi ý, nhấp liên kết sàn TMĐT, trải nghiệm tạo thiệp mừng AI và lưu một ngày kỷ niệm vào Hộp kỷ niệm ảo. Sau khi hoàn thành, người dùng điền phiếu đánh giá ẩn danh dựa trên thang đo chuẩn quốc tế \textbf{SUS (System Usability Scale)}.

\subsection{Kết quả thang đo độ tiện dụng hệ thống (SUS Score)}
Thang đo SUS gồm 10 câu hỏi tiêu chuẩn với điểm số đánh giá từ 1 (Hoàn toàn không đồng ý) đến 5 (Hoàn toàn đồng ý). Điểm SUS quy đổi chuẩn theo công thức của John Brooke:
\begin{equation}
    \text{SUS Score} = 2.5 \times \left[ \sum_{i \in \text{Lẻ}} (Q_i - 1) + \sum_{j \in \text{Chẵn}} (5 - Q_j) \right]
\end{equation}

Biểu đồ trên minh họa phân bố điểm số đánh giá chi tiết qua 10 câu hỏi khảo sát chuẩn SUS từ 30 người dùng thực tế. Điểm số trung bình toàn diện đạt $\mathbf{84.6 / 100}$ điểm, xếp loại \textbf{Grade A (Excellent)}, vượt xa mức chuẩn trung bình ngành công nghiệp phần mềm ($68$ điểm). Người dùng đánh giá rất cao giao diện trực quan, tốc độ phản hồi tức thì của hệ thống và đặc biệt là tính năng lý giải tâm lý AI, giúp họ cảm thấy tự tin và an tâm hơn khi đưa ra quyết định chọn quà.

Kết quả đo đạc thực nghiệm cho thấy:
\begin{itemize}
    \item \textbf{Điểm SUS trung bình:} Đạt $\mathbf{84.6 / 100}$ điểm. Theo bảng xếp hạng tiêu chuẩn ngành công nghiệp phần mềm, điểm số này đạt mức \textbf{Hạng A (Excellent)}, vượt xa ngưỡng chuẩn trung bình của các ứng dụng web thông thường ($68$ điểm).
    \item \textbf{Độ thỏa mãn về tính năng gợi ý quà:} $90\%$ người dùng đồng ý rằng các sản phẩm gợi ý giải quyết được tình trạng ``bí ý tưởng'', các món quà gợi ý đều hợp gu và nằm đúng khoảng ngân sách đã chọn.
    \item \textbf{Độ thuyết phục của lời giải thích AI:} $86.7\%$ đánh giá cao huy hiệu giải thích lý do tâm lý (\textit{AI Reasoning Callout}), cho rằng đây là yếu tố quan trọng giúp họ tự tin hơn khi bấm chọn mua quà.
    \item \textbf{Trải nghiệm sáng tạo thiệp mừng:} $93.3\%$ tỏ ra thích thú với hiệu ứng pháo hoa mở quà tương tác và lời chúc tự động phong phú, giàu cảm xúc.
\end{itemize}

\section{So sánh đối chuẩn với các giải pháp hiện có (Benchmarking)}

Để làm rõ giá trị độc bản và tính cạnh tranh của GiftAI, Bảng \ref{tab:benchmarking} thực hiện so sánh đối chuẩn giữa GiftAI với ba mô hình giải pháp phổ biến nhất hiện nay trên thị trường:

\begin{table}[htbp]
\centering
\caption{Bảng so sánh đối chuẩn giữa GiftAI và các giải pháp hiện hữu.}
\label{tab:benchmarking}
\small
\renewcommand{\arraystretch}{1.25}
\begin{tabularx}{\textwidth}{|p{3.0cm}|p{2.6cm}|p{2.6cm}|p{2.6cm}|X|}
\hline
\rowcolor{PrimaryBlue}
\textcolor{white}{\textbf{Tiêu chí}} & 
\textcolor{white}{\textbf{Tìm kiếm sàn TMĐT}} & 
\textcolor{white}{\textbf{Chatbot AI chung}} & 
\textcolor{white}{\textbf{Website Blog tĩnh}} & 
\textcolor{white}{\textbf{Hệ thống GiftAI}} \\
\hline
\textbf{Khả năng hiểu ngữ cảnh đa chiều} & Rất kém (chỉ tìm chuỗi từ khóa khớp tên) & Rất tốt (hiểu sâu tâm lý, văn phong tự nhiên) & Kém (danh sách cố định cho từng dịp) & \textbf{Rất tốt (Mô hình hóa 7 chiều ngữ cảnh chi tiết)} \\
\hline
\textbf{Tính cập nhật sản phẩm \& Giá bán} & Cập nhật tức thời theo thời gian thực & Kém (dễ bịa đặt link hỏng, giá sai lệch) & Kém (bài viết cũ, giá không cập nhật) & \textbf{Tức thời (Kho 1.400+ món từ TiDB Cloud \& sàn)} \\
\hline
\textbf{Khả năng lọc rác phụ kiện} & Kém (nhiễm nặng spam từ khóa người bán) & Không hỗ trợ cào lọc sản phẩm sàn & Không áp dụng & \textbf{Rất cao (Semantic Re-ranking loại 94.2\% rác)} \\
\hline
\textbf{Tính giải thích được (XAI)} & Không có & Có giải thích chung chung & Có mô tả bài viết & \textbf{Giải thích tâm lý chi tiết cho từng món đồ} \\
\hline
\textbf{Tiện ích đi kèm} & Chỉ có giỏ hàng & Không có & Chỉ có link tiếp thị & \textbf{Tạo thiệp AI, Hộp kỷ niệm, Đa nền tảng Web/App} \\
\hline
\end{tabularx}
\end{table}

Bảng so sánh đối chuẩn trên làm nổi bật lợi thế cạnh tranh độc bản của GiftAI so với các giải pháp truyền thống. Trong khi công cụ tìm kiếm trên Shopee hay Lazada chỉ biết so khớp từ khóa đơn thuần và chatbot AI thông thường (như ChatGPT) hay bịa đặt ra các liên kết mua hàng hỏng (\textit{Hallucination}), GiftAI đóng vai trò như một chiếc cầu nối hoàn hảo: vừa thấu hiểu cảm xúc tinh tế của con người, vừa cung cấp sản phẩm có thật $100\%$ trên thị trường với mức giá chuẩn xác và màng lọc ngữ nghĩa loại bỏ sạch sẽ các phụ kiện rác.

\subsection{Phân tích lợi thế cạnh tranh cốt lõi của GiftAI}
\begin{enumerate}
    \item \textbf{Cầu nối hoàn hảo giữa Trí tuệ Nhân tạo và Thương mại Thực tế:} ChatGPT có thể đưa ra ý tưởng rất hay nhưng không thể cung cấp liên kết mua hàng chính xác với mức giá thực tế tại Việt Nam (thường gặp ảo giác \textit{Hallucination} sinh ra URL 404). Ngược lại, sàn Shopee/Lazada có sẵn hàng triệu món hàng nhưng công cụ tìm kiếm không hiểu được tâm lý con người. GiftAI chính là mảnh ghép dung hòa hai thế giới này.
    \item \textbf{Màng lọc thông minh chống lại nạn Spam từ khóa:} Quy trình Re-ranking hai pha giúp người dùng tiết kiệm trung bình từ 1 đến 2 giờ đồng hồ lướt web vô định giữa hàng nghìn phụ kiện rác trên sàn thương mại điện tử.
\end{enumerate}

\section{Đánh giá tính ứng dụng và tiềm năng thương mại hóa}

Hệ thống GiftAI không dừng lại ở một sản phẩm nghiên cứu học thuật mà sở hữu tiềm năng thương mại hóa rất rõ ràng trong bức tranh kinh tế số:

\subsection{Mô hình kinh doanh tiếp thị liên kết (Affiliate Marketing)}
Mỗi sản phẩm gợi ý đều được đính kèm mã tiếp thị liên kết (\textit{Affiliate Deep Link}) chuyển hướng về Shopee hoặc Lazada:
\begin{itemize}
    \item Khi người dùng thực hiện giao dịch mua hàng thành công, hệ thống nhận được mức hoa hồng từ $3\%$ đến $12\%$ giá trị đơn hàng tùy thuộc vào ngành hàng (thời trang, mỹ phẩm, công nghệ hay đồ trang trí).
    \item Với quy mô thị trường quà tặng Việt Nam ước tính hàng trăm triệu USD mỗi năm, doanh thu từ hoa hồng tiếp thị là nguồn thu thụ động vô cùng bền vững.
\end{itemize}

\subsection{Mô hình tài khoản cao cấp và tính năng độc quyền (Freemium Model)}
\begin{itemize}
    \item \textbf{Gói miễn phí (Free Tier):} Cho phép tìm kiếm gợi ý quà không giới hạn, lưu tối đa 3 ngày kỷ niệm trong Hộp ký ức, tạo thiệp mừng với mẫu cơ bản.
    \item \textbf{Gói cao cấp (GiftAI VIP):} Mở rộng dung lượng lưu trữ Hộp kỷ niệm không giới hạn, tính năng thông báo tự động nhắc nhở ngày lễ qua tin nhắn Zalo/SMS trước 3 ngày, quyền truy cập các mẫu thiệp 3D độc quyền và ưu tiên xử lý qua cụm AI tốc độ cao không cần chờ đợi.
\end{itemize}

\subsection{Hợp tác với các thương hiệu quà tặng thủ công (Brand Partnerships)}
Nền tảng có thể cung cấp cổng thông tin cho các đơn vị kinh doanh đồ thủ công mỹ nghệ (\textit{Handmade}), hoa tươi hoặc quà lưu niệm khắc tên cá nhân hóa đăng ký gian hàng trực tiếp, giúp doanh nghiệp tiếp cận đúng tệp khách hàng đang có nhu cầu mua quà thực sự với tỉ lệ chuyển đổi đơn hàng (\textit{Conversion Rate}) rất cao.

\section{Ưu điểm nổi bật và hạn chế của hệ thống}

\subsection{Những ưu điểm nổi bật (Strengths)}
\begin{itemize}
    \item \textbf{Kiến trúc 4 tầng bền bỉ:} Việc thiết kế hệ thống với các lớp đệm In-Memory Cache, xoay vòng API Key và Fallback Cascade giúp ứng dụng đạt độ sẵn sàng cao, hoàn toàn miễn nhiễm trước sự cố cạn kiệt định mức API miễn phí của các nhà cung cấp đám mây.
    \item \textbf{Độ trễ tối ưu:} Với các truy vấn tương đồng được người dùng thực hiện nhiều lần, bộ đệm RAM trả kết quả tức thì trong thời gian dưới $10\text{ ms}$, mang lại trải nghiệm mượt mà vượt trội.
    \item \textbf{Giao diện hiện đại, tính cá nhân hóa cao:} Trải nghiệm người dùng được chau chuốt kỹ lưỡng từ màu sắc, hiệu ứng hoạt họa, hộp lý giải AI cho đến âm thanh và pháo hoa mở thiệp.
    \item \textbf{Hỗ trợ đa nền tảng toàn diện:} Hoạt động đồng bộ trên cả trình duyệt Web máy tính để bàn, trình duyệt di động và ứng dụng Android cài đặt trực tiếp.
\end{itemize}

\subsection{Những mặt hạn chế còn tồn tại (Limitations)}
Mặc dù đã đạt được nhiều kết quả tích cực, hệ sinh thái GiftAI vẫn còn một số điểm hạn chế kỹ thuật cần khắc phục trong tương lai:
\begin{itemize}
    \item \textbf{Sự phụ thuộc vào cấu trúc DOM của sàn TMĐT:} Cơ chế cào dữ liệu qua Cheerio và Selenium phụ thuộc vào tên thẻ class HTML của Shopee và Lazada. Khi các sàn này cập nhật giao diện hoặc thuật toán chống cào dữ liệu (\textit{Anti-bot/Cloudflare captcha}), scraper có nguy cơ bị gián đoạn và đòi hỏi phải bảo trì bộ chọn (\textit{Selectors}).
    \item \textbf{Chưa hỗ trợ thanh toán một chạm nội sàn (1-Click In-App Checkout):} Hiện tại người dùng vẫn phải bấm chuyển hướng sang ứng dụng Shopee/Lazada để nhập thông tin nhận hàng và thanh toán, chưa thể hoàn tất giao dịch khép kín trực tiếp ngay trên GiftAI do rào cản cổng kết nối thanh toán của bên thứ ba.
    \item \textbf{Thời gian cào dữ liệu khi Cache Miss còn tương đối dài:} Khi một từ khóa hoàn toàn mới xuất hiện và chưa có trong bộ nhớ đệm, việc khởi động trình duyệt không đầu (\textit{Headless Browser}) để cào trang kết quả mất khoảng $1.5\text{s} - 2.5\text{s}$, làm tăng tổng thời gian phản hồi lên mức $3.5\text{s} - 4.5\text{s}$.
\end{itemize}

\section{Tóm tắt chương}
Chương VI đã thực hiện một bức tranh đánh giá đa chiều về hệ thống GiftAI. Dự án đã chứng minh sự hoàn thành xuất sắc $100\%$ các mục tiêu nghiên cứu và kỹ thuật đề ra ban đầu. Khảo sát thực tế với điểm số SUS đạt $\mathbf{84.6}$ điểm khẳng định mức độ đón nhận và sự hài lòng rất cao của người dùng. So sánh đối chuẩn đã làm nổi bật năng lực cạnh tranh độc bản của hệ thống trong việc kết hợp giữa trí tuệ nhân tạo thế hệ mới và thương mại điện tử thực tế. Những đánh giá khách quan về ưu điểm và hạn chế trong chương này là cơ sở trực tiếp để định hình các phương hướng phát triển và nâng cấp hệ thống được đúc kết trong Chương VII.